Welcome to The Third Arabic Natural Language Processing Conference (ArabicNLP 2025), co-located with EMNLP 2025 in Suzhou, China. ArabicNLP 2029 is the tenth edition of the WANLP/ArabicNLP meeting series, which has developed a growing reputation as a high quality venue for researchers and engineers working on Arabic NLP, where they share and discuss their ongoing work. The first in the WANLP series was held in Doha, Qatar (EMNLP 2014), followed by Beijing, China (ACL 2015), Valencia, Spain (EACL 2017), Florence, Italy (ACL 2019), online with COLING 2020, online with EACL 2021, a hybrid event in Abu Dhabi, UAE (EMNLP 2022), and then finally in-person events in Singapore with EMNLP 2023 and Bangkok, Thailand with ACL 2024.

For this year’s edition of ArabicNLP, we received a total of 95 correct main conference submissions and accepted 39 papers, which brings us to an acceptance rate of 41.0\%, the lowest to date in the WANLP/ArabicNLP series. All papers submitted to the conference were reviewed by at least three reviewers each.

ArabicNLP 2025 included eleven shared tasks with 138 papers (11 overview papers and 127 system descriptions) in three main tracks:

\textbf{Track 1: Speech and Multimodal Processing }
\begin{enumerate}
\item  ImageEval Arabic Image Captioning with X papers
Iqra'Eval: A Shared Task on Qur'anic Pronunciation Assessment with X papers
\item NADI 2025: Multidialectal Arabic Speech Processing with X papers
\item MAHED 2025: Multimodal Detection of Hope and Hate Emotions in Arabic Content
\end{enumerate}

\textbf{Track 2: Text Quality and Generation Assessment}
\begin{enumerate}
\item AraGenEval: Arabic Authorship Style Transfer and AI Generated Text Detection
\item TAQEEM 2025: The First Task for Arabic Quality Evaluation of Essays in Multi-dimensions
\item BAREC 2025: Arabic Readability Assessment Shared Task
\item AraHealthQA 2025: Comprehensive Arabic Health Question Answering Shared Task
\end{enumerate}

\textbf{Track 3: Cultural and Ethical Evaluation of LLMs for Arabic}
\begin{enumerate}
\item IslamicEval: Capturing LLMs Hallucination in Islamic Content
\item PalmX 2025: The First Shared Task on Benchmarking LLMs on Arabic Culture
\item QIAS 2025: Q\&A in Islamic Studies Assessment Shared Task
\end{enumerate}

ArabicNLP 2025 also includes two invited talks by Houda Bouamor, entitled “Beyond Resources: Building an Arabic NLP Ecosystem Rooted in Representation, Collaboration, and Responsibility” and Areeb Alowisheq, entitled “From Benchmarks to the Real-World Impact: Arabic LLMs in Production”. We were able to secure sponsorship funding from different institutions, Humain, Google, CAMeL Lab (NYU-AD), and Clinical AI Lab (NYU-AD). We used the sponsorship funds to support student registrations. We thank all our sponsors for their generous support and their help in building up the Arabic NLP community. Finally, we extend our gratitude to everyone who submitted a paper to the conference, and to the Program Committee members for their diligent efforts in providing reviews within a very tight time frame.

Kareem Darwish, General Chair, on behalf of the conference organizers.

Website of the conference: https://arabicnlp2025.sigarab.org